\documentclass[10pt,a4paper]{amsart}
\usepackage[margin=19mm]{geometry}
\usepackage{amsmath,amssymb,mathtools}
\usepackage[scr=rsfs,cal=euler]{mathalfa}
\usepackage{xfrac}
\usepackage[unicode,hidelinks,bookmarks=false]{hyperref}
\newcommand{\ie}{i.e.\ }
\newcommand{\cf}{cf.\ }
\newcommand{\resp}{respectively\ }
\newcommand{\Subsection}{\subsection}
\providecommand{\nobreakdash}{\nobreak\hbox{-}}
\setlength{\emergencystretch}{2em}
\multlinegap0pt
\usepackage{mdframed}
\usepackage{mdframed}
\usepackage{mdframed}
\usepackage{mdframed}
\usepackage{bm}
\usepackage{amssymb}
\usepackage{xcolor}
\usepackage{enumerate}
\usepackage{mdframed}
\newtheorem{prop}{Proposition}
\title{On Duality, Legendre Bundles and Deformations}
\author{N.C. Combe, P.G. Combe, H.K. Nencka}
\date{Source: arXiv:2604.04865v1 (CC BY 4.0)}
\begin{document}
\maketitle

\noindent\emph{Source and licence.} On Duality, Legendre Bundles and Deformations. Version arXiv:2604.04865v1 (CC BY 4.0), distributed by arXiv under the Creative Commons Attribution 4.0 International licence, http://creativecommons.org/licenses/by/4.0/, as recorded in the arXiv licence metadata for this exact version and shown on the abstract page https://arxiv.org/abs/2604.04865v1. Full licence notice: \texttt{licenses/arxiv-2604.04865-CC-BY-4.0.txt}.\par\smallskip

\noindent\textit{Editorial scope.} Counted unit: the article's prop P:ParaKahler (source0.tex lines 455--457). The article's complete original proof of it is delivered with it (source0.tex lines 458--490, one \texttt{proof} environment of 33 lines). No mathematics is written, repaired or re-derived: the statement and the proof are the article's own lines, in the article's own notation, and no retained line is substituted or re-flowed.  No further source-local context is needed: every cross-reference of every retained line resolves inside the two delivered spans.  Nothing was omitted from any retained span, so the package carries no omission record.  No retained line contains a citation, so the wrapper carries no bibliography and no bibliography entry.  No retained line contains a figure environment, an image-inclusion command or drawing code, so no image is delivered and the package declares no asset dependency.  Editorial formatting only, in addition: an AMS \texttt{amsart} preamble that restates the article's own declarations and macros used by the retained lines, namely the environments \texttt{prop} on separate counters, together with the standard packages those lines need.  The retained environments print in this document as Proposition 1 in order of appearance, whereas the article prints the same items with the numbering of its own full document.  The complete native source of the article as distributed in the arXiv e-print for this exact version is bundled unchanged as \texttt{sources/197915/source0.tex}.
\begin{prop}\label{P:ParaKahler}
The (Legendre) bundle $(H,J,\boldsymbol{\omega}^S)$ is a para‑Kähler vector bundle over $B$. 
\end{prop}

    \begin{proof}
Given the decomposition $H = H^+ \oplus H^-$, define $X = X^+ \oplus \alpha$ and $Y = Y^+ \oplus \beta,$ where $X^+, Y^+ \in \Gamma(H^+)$ and $\alpha, \beta \in \Gamma(H^-)$. Then,
by definition of $J$ we obtain \[
JX = X^+ \oplus (-\alpha).
\]
Using the pairing
\[
\langle\!\langle X^+, \beta \rangle\!\rangle = \beta(X^+)
\quad \text{and} \quad
\langle\!\langle \alpha, Y^+ \rangle\!\rangle = \alpha(Y^+),
\]
we define  $$\boldsymbol{\omega}^S(X,Y)=\langle\!\langle 
 J(X \oplus \alpha),\, Y \oplus \beta \rangle\!\rangle
= \langle\!\langle X^+ \oplus (-\alpha),\, Y^+ \oplus \beta \rangle\!\rangle$$ \[
= \langle\!\langle X^+, Y^+ \rangle\!\rangle
+ \langle\!\langle X^+, \beta \rangle\!\rangle
+ \langle\!\langle -\alpha, Y^+ \rangle\!\rangle
+ \langle\!\langle -\alpha, \beta \rangle\!\rangle
\]
\[
= 0 + \beta(X^+) + \bigl(-\alpha(Y^+)\bigr) + 0
= \beta(X^+) - \alpha(Y^+).
\] which is the canonical symplectic form on $T^*B$ (when identifying $H^+ \cong TB$, $H^- \cong T^*B$).



The natural connection on $H$ is the direct sum $\nabla = \nabla^+ \oplus \nabla^-$, defined by
\[
\nabla(X^+ \oplus \alpha) = (\nabla^+ X^+) \oplus (\nabla^- \alpha).
\]
This connection preserves the splitting, hence it commutes with $J$ (since $J$ is constant on each factor). Moreover, because the pairing $\langle\!\langle \cdot, \cdot \rangle\!\rangle$ is flat with respect to $\nabla^+ \oplus \nabla^-$ (by the duality condition), the symplectic form $\boldsymbol{\omega}^S$ is also flat. Thus $(H, J, \boldsymbol{\omega}^S, \nabla)$ is a flat para-Kähler vector bundle.

    \end{proof}

\end{document}
